\documentclass[aspectratio=169,11pt]{beamer}

% Pakete
\usepackage[utf8]{inputenc}
\usepackage[T1]{fontenc}
\usepackage[ngerman]{babel}
\usepackage{amsmath,amssymb}
\usepackage{tikz}
\usetikzlibrary{arrows.meta,positioning}
\usepackage{siunitx}
\sisetup{locale=DE, per-mode=fraction}

% Beamer-Theme: schlicht
\usetheme{default}
\usecolortheme{default}
\setbeamertemplate{navigation symbols}{}
\setbeamertemplate{footline}[frame number]

% Fachfarbe Physik
\definecolor{physik}{HTML}{2c5aa0}
\definecolor{physiklight}{HTML}{e3f2fd}
\setbeamercolor{structure}{fg=physik}
\setbeamercolor{frametitle}{fg=physik}
\setbeamercolor{block title}{bg=physik,fg=white}
\setbeamercolor{block body}{bg=physiklight}

% Titel
\title{Der Photoeffekt}
\subtitle{Einsteinsche Gerade und Bestimmung von \textit{h}}
\author{Physik 12 GK}
\date{09.01.2026}

\begin{document}

% ============================================================
% TITELFOLIE
% ============================================================
\begin{frame}
\titlepage
\end{frame}

% ============================================================
% FOLIE 1: WIEDERHOLUNG (Einstieg)
% ============================================================
\begin{frame}{Wiederholung: Das Problem}

\begin{columns}
\column{0.6\textwidth}
\textbf{Hallwachs-Versuch:}
\begin{itemize}
    \item UV-Licht löst Elektronen aus Zink
    \item Sichtbares Licht: \textbf{keine Wirkung}
    \item Auch nicht bei hoher Intensität!
\end{itemize}

\vspace{1em}
\textbf{Widerspruch zum Wellenmodell:}
\begin{itemize}
    \item Mehr Intensität = mehr Energie?
    \item Sollte irgendwann Elektronen lösen...
    \item \textbf{Aber: Frequenz entscheidet!}
\end{itemize}

\column{0.35\textwidth}
\begin{tikzpicture}[scale=0.8]
    % Zinkplatte
    \draw[fill=gray!30] (0,0) rectangle (2,3);
    \node at (1,1.5) {Zn};
    \node at (1,0.5) {\small negativ};

    % UV-Licht
    \draw[->,thick,violet] (3.5,2.5) -- (2.2,2.5);
    \draw[->,thick,violet] (3.5,2) -- (2.2,2);
    \node[right,violet] at (3.5,2.25) {UV};

    % Elektronen
    \draw[->,thick,blue] (2.2,2.8) -- (3,3.2);
    \draw[->,thick,blue] (2.2,1.7) -- (3,1.3);
    \node[right,blue] at (3,2) {\small $e^-$};
\end{tikzpicture}
\end{columns}

\end{frame}

% ============================================================
% FOLIE 2: EINSTEINS HYPOTHESE (Input + Formel)
% ============================================================
\begin{frame}{Einsteins Lichtquanten-Hypothese (1905)}

\begin{block}{Die revolutionäre Idee}
Licht besteht aus \textbf{Energiepaketen} -- \textbf{Photonen}
\end{block}

\vspace{0.5em}

\begin{center}
\huge
$\boxed{E_{\text{Photon}} = h \cdot f}$
\end{center}

\vspace{0.5em}

\begin{columns}
\column{0.45\textwidth}
\begin{tabular}{ll}
$h$ & = \SI{6,63e-34}{J.s} \\
    & = \SI{4,14e-15}{eV.s} \\
$f$ & = Frequenz des Lichts \\
\end{tabular}

\vspace{0.5em}
\textbf{Beispiele:}\\
{\small
Rot: $E = \SI{2,1}{eV}$\\
Blau: $E = \SI{3,1}{eV}$\\
UV: $E = \SI{4,1}{eV}$
}

\column{0.52\textwidth}
\begin{alertblock}{Kernidee: Energie ist gequantelt!}
\begin{itemize}
    \item Ein Photon = \textbf{ein} Paket = $h \cdot f$
    \item \textbf{Unteilbar} (wie eine Münze)
    \item Mehr Intensität = mehr Photonen\\
          \textbf{nicht} mehr Energie pro Photon!
\end{itemize}
\end{alertblock}
\end{columns}

\vspace{0.3em}
\begin{center}
\small\textit{Nobelpreis 1921 für diese Erklärung!}
\end{center}

\end{frame}

% ============================================================
% FOLIE 3: GEGENFELDMETHODE (Gegenfeld-Erklärung)
% ============================================================
\begin{frame}{Die Gegenfeldmethode}

\textbf{Frage:} Wie messen wir die kinetische Energie der Elektronen?

\vspace{0.5em}

\begin{columns}
\column{0.55\textwidth}
\textbf{Prinzip:}
\begin{enumerate}
    \item Licht löst Elektronen aus (Kathode)
    \item Elektrisches \textbf{Gegenfeld} bremst sie
    \item Spannung erhöhen bis Strom = 0
    \item Diese \textbf{Grenzspannung} $U_G$ misst $E_{\text{kin}}$
\end{enumerate}

\vspace{0.5em}

\begin{block}{Energieerhaltung}
\centering
$E_{\text{kin,max}} = e \cdot U_G$
\end{block}

\column{0.4\textwidth}
\begin{tikzpicture}[scale=0.7]
    % Kathode
    \draw[fill=gray!30] (0,0) rectangle (0.5,3);
    \node[rotate=90] at (0.25,1.5) {\small Kathode};

    % Anode
    \draw[fill=gray!30] (4,0) rectangle (4.5,3);
    \node[rotate=90] at (4.25,1.5) {\small Anode};

    % Licht
    \draw[->,thick,violet] (-1.5,2) -- (-0.2,2);
    \draw[->,thick,violet] (-1.5,1.5) -- (-0.2,1.5);
    \node[left,violet] at (-1.5,1.75) {$h \cdot f$};

    % Elektronen
    \draw[->,thick,blue] (0.7,2) -- (2,2);
    \draw[->,thick,blue] (0.7,1) -- (1.5,1);
    \node[above,blue] at (1.3,2) {\small $e^-$};

    % Gegenfeld (E zeigt von + nach -, also von Anode nach Kathode)
    \draw[->,thick,red] (3.8,1.5) -- (0.7,1.5);
    \node[above,red] at (2.2,1.5) {\small $\vec{E}$};

    % Spannung
    \draw (0.25,-0.5) -- (0.25,-1);
    \draw (4.25,-0.5) -- (4.25,-1);
    \draw (0.25,-0.75) -- (4.25,-0.75);
    \node[below] at (2.25,-0.75) {$U_G$};
    \node at (1,-0.75) {\small $-$};
    \node at (3.5,-0.75) {\small $+$};
\end{tikzpicture}
\end{columns}

\vspace{0.5em}
\textbf{Praktisch:} Bei $U_G$ in Volt ist $E_{\text{kin}}$ numerisch gleich in \textbf{eV}!

\end{frame}

% ============================================================
% FOLIE 4: MESSUNG (Partnerarbeit mit Simulation)
% ============================================================
\begin{frame}{Jetzt seid ihr dran!}

\begin{center}
\Large
\textbf{PhET-Simulation: Photoeffekt}

\vspace{1em}

\normalsize
\url{phet.colorado.edu/de/simulations/photoelectric}

\vspace{2em}

\textbf{Aufgabe:}\\
Miss $U_G$ für verschiedene Frequenzen (Natrium)\\
und zeichne die Einsteinsche Gerade!

\vspace{2em}

$\rightarrow$ \textbf{AB-02, Aufgaben 1--3}
\end{center}

\end{frame}

% ============================================================
% FOLIE 5: DIE EINSTEINSCHE GERADE (Ableitung)
% ============================================================
\begin{frame}{Die Einsteinsche Gerade}

\begin{columns}
\column{0.55\textwidth}

\begin{tikzpicture}[scale=0.8]
    % Vereinfachtes Diagramm für Projektion
    % Natrium: W_A = 2,3 eV, f_G = 5,56·10^14 Hz

    % Achsen
    \draw[->] (0,0) -- (8.5,0) node[right] {$f$};
    \draw[->] (0,-2.8) -- (0,1.5) node[above] {$E_{\text{kin}}$};

    % Achsenbeschriftung (nur wichtige Werte)
    \foreach \x in {2,4,6,8} {
        \draw (\x,0.08) -- (\x,-0.08);
    }
    \node[below] at (6,0) {\scriptsize $6$};
    \node[below] at (8,0) {\scriptsize $8 \cdot 10^{14}$ Hz};
    \draw (0.08,1) -- (-0.08,1) node[left] {\scriptsize 1 eV};
    \node[below left] at (0,0) {\scriptsize 0};

    % Gerade: E = h·f - W_A (gestrichelt bis f_G, dann durchgezogen)
    \draw[dashed,thick,gray] (0,-2.3) -- (5.56,0);
    \draw[very thick,physik] (5.56,0) -- (8,1.01);

    % Messpunkte
    \fill[physik] (6.0,0.18) circle (3pt);
    \fill[physik] (7.0,0.60) circle (3pt);
    \fill[physik] (8.0,1.01) circle (3pt);

    % Grenzfrequenz f_G
    \fill[red] (5.56,0) circle (3pt);
    \node[below,red] at (5.56,-0.15) {\small $f_G$};

    % y-Achsenabschnitt -W_A
    \fill[orange] (0,-2.3) circle (3pt);
    \draw[dashed,orange] (0,-2.3) -- (-0.3,-2.3);
    \node[left,orange] at (-0.3,-2.3) {\small $-W_A$};

    % Steigung = h (einfache Beschriftung)
    \node[physik,rotate=24] at (7,0.9) {\small Steigung $= h$};
\end{tikzpicture}

\column{0.40\textwidth}
\textbf{Interpretation (Na):}

\vspace{0.3em}

\begin{tabular}{ll}
\textbf{Steigung} & $= h$ \\
& {\scriptsize (universell!)} \\[0.3em]
\textbf{$x$-Abschnitt} & $= f_G$ \\
& {\scriptsize $= \SI{5,56e14}{Hz}$} \\[0.3em]
\textbf{$y$-Abschnitt} & $= -W_A$ \\
& {\scriptsize $= \SI{-2,3}{eV}$} \\
\end{tabular}

\vspace{0.8em}

\begin{block}{Photogleichung}
\centering
$E_{\text{kin}} = h \cdot f - W_A$
\end{block}

\vspace{0.3em}
{\small Gestrichelt: Extrapolation\\(dort kein Photoeffekt!)}

\end{columns}

\end{frame}

% ============================================================
% FOLIE 6: ZUSAMMENFASSUNG (Sicherung)
% ============================================================
\begin{frame}{Zusammenfassung: Die Photogleichung}

\begin{center}
\Large
$\boxed{E_{\text{kin,max}} = h \cdot f - W_A}$
\end{center}

\vspace{1em}

\begin{columns}
\column{0.5\textwidth}
\textbf{Größen:}
\begin{tabular}{ll}
$E_{\text{kin}}$ & kinetische Energie \\
$h$ & \SI{4,14e-15}{eV.s} \\
$f$ & Lichtfrequenz \\
$W_A$ & Austrittsarbeit \\
\end{tabular}

\column{0.5\textwidth}
\textbf{Abgeleitete Größen:}
\begin{tabular}{ll}
$f_G$ & $= \dfrac{W_A}{h}$ \\[0.8em]
$U_G$ & $= \dfrac{E_{\text{kin}}}{e}$ \\
\end{tabular}
\end{columns}

\vspace{1em}

\begin{block}{Bedeutung}
\begin{itemize}
    \item Unter $f_G$: \textbf{kein Photoeffekt} (egal wie hell!)
    \item Über $f_G$: $E_{\text{kin}}$ steigt \textbf{linear} mit $f$
    \item Aus der Steigung: \textbf{$h$ bestimmbar}!
\end{itemize}
\end{block}

\end{frame}

\end{document}
